% Chapter 2

\chapter{相关理论与关键技术基础}

\section{扩散模型基本原理}

\subsection{前向扩散过程}

扩散模型是一类通过逐步加噪与逐步去噪完成数据建模的生成模型。Ho等人在Denoising Diffusion Probabilistic Models\cite{ref1}中将生成过程描述为两个方向相反的马尔可夫链：在前向扩散过程中，模型按照预先设定的噪声调度策略，向真实样本逐步注入高斯噪声，使原始数据从清晰图像逐渐演化为近似各向同性高斯分布；在反向过程中，再由神经网络学习如何从噪声中逐步恢复目标样本。与依赖对抗训练的传统GAN相比，扩散模型训练更稳定、模式覆盖能力更强，因此逐渐成为当前图像与视频生成任务的主流框架。

在数学形式上，设原始样本为$x_0$，扩散模型通过一系列时间步$t=1,2,\ldots,T$将样本逐步变换为$x_t$。每一步前向扩散都可以表示为在上一时刻样本上加入少量高斯噪声，因此随着时间步增加，样本中的语义结构与细节信息会被持续削弱，最终得到近似纯噪声的$x_T$。这一过程本身不需要学习参数，它的作用在于把复杂的数据分布转化为更容易处理的噪声分布，从而为反向生成建立统一的概率路径。对于图像生成任务而言，前向扩散可以看作是一个从有意义图像到无结构噪声的连续退化过程。

\subsection{反向去噪过程与训练目标}

扩散模型真正需要学习的是反向去噪过程，即在给定当前噪声样本$x_t$的条件下，预测其在上一时刻$x_{t-1}$的分布，逐步将噪声恢复为符合目标数据分布的样本。由于直接建模完整的反向分布较为困难，DDPM\cite{ref1}将其近似为高斯分布，并利用神经网络预测噪声项或等价变量，从而实现逐步去噪。在实际训练中，通常将目标简化为最小化模型预测噪声与真实噪声之间的均方误差，这种训练方式既保留了理论上的变分推断基础，也显著提高了实现效率。因此，在Stable Diffusion等后续模型中，扩散模型通常表现为一个在不同时间步估计噪声的条件生成器。

在文本条件生成任务中，扩散模型还需要引入条件信息来控制生成方向。Classifier-Free Guidance\cite{ref31}进一步提出，无需额外训练分类器，就可以通过同时学习有条件与无条件两个分支，在采样阶段对二者的预测结果进行线性组合，从而在样本多样性和语义一致性之间取得平衡。该方法已经成为现代文生图模型中的标准采样策略。对于本研究而言，提示词、负向提示词以及guidance scale等参数本质上都与这一引导机制密切相关，它决定了生成结果是更强调文本约束，还是保留更多底模的自由生成能力。

扩散模型中前向加噪与反向去噪之间的对应关系如图\ref{fig:ch2-diffusion-process}所示。

\begin{figure}[htbp]
  \centering
  \includegraphics[width=0.92\textwidth]{chapter2_figure1_1.png}
  \caption{扩散模型的前向与反向扩散过程}
  \label{fig:ch2-diffusion-process}
\end{figure}

\section{潜空间扩散模型与Stable Diffusion}

\subsection{潜空间表示与VAE压缩机制}

早期扩散模型大多直接在像素空间中进行加噪与去噪，这种方式虽然能够取得较高的生成质量，但在高分辨率图像场景下会带来很高的显存和计算开销。Rombach等人在Latent Diffusion Models\cite{ref4}中提出，将扩散过程从像素空间迁移到低维潜空间中进行，即先使用自编码器将图像压缩为更紧凑的潜变量表示，再在潜变量上完成扩散建模，最后通过解码器重建高分辨率图像。这一思路使扩散模型在细节保真和训练成本之间取得了较好的平衡，也是Stable Diffusion系列模型能够在有限算力条件下广泛应用的重要原因。

在潜空间扩散框架中，变分自编码器承担感知压缩和图像重建的任务。编码器负责将原始图像映射到低维潜变量$z$，解码器则将生成得到的潜变量重新还原为可视图像\cite{ref4}。由于潜空间维度远小于原始像素空间，模型在训练时不必直接处理冗余像素细节，而是集中学习更高层次的语义结构与纹理模式。对本研究而言，使用Stable Diffusion v1.5作为底模，本质上是在继承LDM架构优势的基础上开展后续风格学习，因此理解潜空间建模机制，对于分析分辨率设置、显存占用与训练速度之间的关系具有重要意义。

\subsection{U-Net去噪网络与交叉注意力条件注入}

在Stable Diffusion中，负责执行去噪任务的核心网络通常是U-Net。U-Net采用编码器—解码器对称结构，通过下采样逐步提取全局语义信息，再通过上采样逐步恢复空间分辨率，并借助跳跃连接将浅层细节与深层语义进行融合。扩散模型在每一个时间步都会将当前潜变量、时间步嵌入以及文本条件共同输入U-Net，使网络能够在不同噪声强度下学习如何恢复图像结构\cite{ref4}。这种多尺度特征融合能力使U-Net非常适合承担扩散去噪器的角色，也使其成为后续LoRA插入的主要位置。

为了让文本提示词真正影响生成结果，LDM\cite{ref4}在U-Net中引入了交叉注意力机制。文本编码器首先将提示词转换为一组语义嵌入向量，随后这些向量通过交叉注意力与图像潜变量特征发生交互，从而把语言信息注入图像生成过程。这一设计使模型不再只是无条件地从噪声恢复图像，而是可以在每个去噪阶段持续参考文本语义，逐步生成更符合提示词要求的结果。本研究中的吉卜力风格LoRA训练，本质上也是围绕这一条件生成框架展开：一方面保持底模原有的文本到图像映射能力，另一方面通过少量可训练参数改变注意力层对特定风格特征的响应。

\section{LoRA参数高效微调技术}

\subsection{LoRA的基本原理}

LoRA是一种典型的参数高效微调方法，其核心思想是：在不直接更新原始大模型权重的前提下，用两个低秩矩阵近似表示权重更新量。Hu等人\cite{ref3}指出，许多大模型在下游任务适配时，真正有效的权重变化往往具有较低的内在秩，因此没有必要对全部参数进行全量更新。基于这一认识，LoRA将某一层的权重更新写为$\Delta W = BA$，其中$A$和$B$的秩远小于原始权重矩阵的维度，训练时只优化这两个小矩阵，而冻结底模原有参数。这样既能显著减少训练参数量，也不会改变原始模型的主体结构。

LoRA与常规全参数微调在参数更新范围和训练代价上的差异，可由图\ref{fig:ch2-lora-vs-finetune}进一步说明。

\begin{figure}[htbp]
  \centering
  \includegraphics[width=0.92\textwidth]{chapter2_figure2_1.png}
  \caption{LoRA与常规微调的原理对比}
  \label{fig:ch2-lora-vs-finetune}
\end{figure}

LoRA的优势主要体现在三个方面。首先，它大幅降低了显存需求与训练成本，使大模型在消费级显卡上的定制训练成为可能；其次，LoRA训练得到的适配器文件体积较小，便于存储、切换和复用；最后，由于原始底模权重保持不变，多个不同风格或任务的LoRA可以在同一底模上独立训练并灵活加载\cite{ref3}。对于毕业设计这类需要在有限周期内完成多轮实验验证的任务而言，LoRA具有较高的实用性。它使研究者能够把更多精力放在数据集选择、提示词设计、训练步数控制和风格效果分析上，而不必长期受制于全参数微调带来的高昂资源成本。

\subsection{LoRA在扩散模型中的应用方式}

虽然LoRA最初提出于大语言模型场景，但其思想同样适用于扩散模型中的线性映射层，尤其是自注意力层和交叉注意力层中的投影矩阵。对于Stable Diffusion这类模型而言，风格、主体特征以及文本响应能力在很大程度上都通过注意力机制进行组织，因此将LoRA注入到query、key、value或输出投影层，能够以较小代价改变模型对视觉风格特征和语义条件的表达方式。目前社区中广泛采用的SD-LoRA方案\cite{ref3,ref4}，正是基于这一技术路线形成的。

在本研究的实现中，LoRA主要承担特定风格适配器的角色。训练时冻结底模的VAE、文本编码器和大部分U-Net参数，仅对注入的低秩矩阵进行更新，使模型逐步学习目标风格数据中的颜色分布、画面纹理、场景构图与笔触特征。由于训练参数较少，研究者可以通过调节rank、学习率、训练步数等超参数，快速比较不同实验配置对生成结果的影响。进一步地，当LoRA与视频生成或结构控制模块结合时，它还可以作为一种可插拔的风格组件，与运动模块、ControlNet等机制形成协同，这也是本研究后续视频阶段实验的重要理论基础。

\section{视觉—语言预训练与数据标注技术}

\subsection{CLIP模型}

CLIP\cite{ref23}是当前视觉—语言预训练领域的代表性工作之一。Radford等人通过对大规模图文对数据进行对比学习，使图像编码器和文本编码器在同一嵌入空间中完成对齐，从而让模型能够学习图像内容与自然语言描述之间的对应关系。与仅面向固定类别的传统图像分类模型不同，CLIP利用自然语言作为监督信号，使视觉模型具备更强的开放词汇理解能力和跨任务迁移能力。这种能力对于生成模型尤为重要，因为文生图任务本身就依赖文本条件与视觉语义之间的有效对应。

在扩散模型生态中，CLIP的作用主要体现在两个方面。其一，它可以作为文本—图像一致性评估工具，用于衡量生成图像是否与给定提示词相匹配；其二，它也可以参与数据筛选、图文检索和风格样本组织，帮助提升训练数据的语义质量\cite{ref23}。对于本研究而言，CLIP相关思想有助于理解高质量caption为何对LoRA风格训练如此重要：如果文本描述无法准确概括样本中的场景、主体和艺术风格特征，那么模型在学习时就难以建立稳定的文本—图像映射关系，最终会影响生成阶段的可控性与复现性。

\subsection{BLIP与BLIP-2自动标注机制}

仅依靠原始互联网图文对进行训练，往往会受到数据噪声较多、标注不精确等问题影响。BLIP\cite{ref32}提出通过“生成伪标题+过滤噪声样本”的方式，对大规模图文数据进行自举式改进，使模型同时具备视觉理解与文本生成能力。这一机制的价值在于，它不仅能够完成图文检索、视觉问答等理解任务，还能生成图像描述，因此非常适合作为数据预处理阶段的自动caption工具。对于风格学习任务而言，BLIP能够帮助研究者为样本补充更完整的场景、主体和风格描述，从而提高训练样本的文本质量。

BLIP-2\cite{ref33}进一步提出，在冻结视觉编码器与大语言模型的前提下，通过轻量级Querying Transformer桥接视觉特征与语言生成过程，从而以更低的训练代价获得较强的图文理解与生成能力。相比端到端训练大规模视觉语言模型，BLIP-2的实现方式更高效，也更适合在实际工程中作为自动标注模块使用。本研究前期的数据构建思路正与此类模型高度契合：通过自动生成caption，再结合人工筛查和风格词修正，使训练数据不仅保留客观内容信息，还能突出“手绘感”“动画电影风格”“暖色光影”等风格特征，为LoRA学习提供更明确的监督信号。

\section{视频生成与时序建模技术}

\subsection{基于扩散模型的视频生成}

视频生成可以看作是图像生成向时间维度的自然扩展，但其难点并不仅仅在于处理更多帧，而在于同时保证单帧质量和跨帧一致性。Video Diffusion Models\cite{ref34}指出，视频扩散模型可以视为标准图像扩散架构在时空维度上的推广，通过在训练中联合利用图像与视频数据，不仅能够提升生成视频的清晰度与稳定性，还能降低小批量训练时梯度估计的方差，加快优化过程。这说明视频生成并不是完全脱离图像生成而独立存在的任务，而是与图像扩散模型共享大量基础结构与学习机制。

在具体实现层面，视频扩散模型通常需要在网络中显式建模时间维度信息，例如通过三维卷积、时空注意力或帧间特征交互模块，使模型能够感知相邻帧之间的运动变化与语义连续性\cite{ref34}。如果仅把每一帧独立当作图像生成，虽然可以获得较好的单帧质量，但往往会在视频播放时出现闪烁、主体漂移和动作不连贯等问题。因此，时序一致性建模是视频生成区别于图像生成的关键所在。本研究之所以在图像风格LoRA之外继续关注视频阶段，就是希望在已有风格学习结果的基础上，进一步分析风格表达与时间稳定性之间的协调关系。

\subsection{AnimateDiff运动模块}

AnimateDiff\cite{ref12}是近年来图像扩散模型向视频生成扩展的代表性方法之一。该方法的核心贡献在于提出了一种可插拔的运动模块，可以在不对每一个个性化图像模型分别进行视频微调的前提下，直接把已有的文本到图像模型激活为动画生成器。其基本思想是：把外观和风格能力交给原有的图像扩散模型处理，把时间动态和运动先验交给单独训练的motion module负责。这样一来，研究者就可以在保留LoRA、DreamBooth等个性化外观能力的同时，引入统一的时间运动建模能力。

AnimateDiff还提出了MotionLoRA，用于以较小成本适配新的镜头类型或运动模式\cite{ref9}。这一思路对本研究具有直接启发意义。因为本研究的重点并不是从零开始训练一个大规模视频基础模型，而是在Stable Diffusion + LoRA的图像风格学习框架上，进一步扩展到短视频生成验证。在这种情况下，采用底模负责外观、附加模块负责运动的解耦思路，比全量训练视频模型更符合时间和算力上的现实条件，也更容易形成可复现实验闭环。

\section{可控生成与结构约束技术}

\subsection{ControlNet条件控制机制}

当生成任务不仅要求接近某种风格，还要求保持特定结构、姿态或布局时，单纯依靠文本提示词和风格LoRA往往是不够的。ControlNet\cite{ref14}针对这一问题提出，在冻结原有扩散模型主干的基础上，增加一组用于接收结构条件的控制分支，并通过零初始化卷积逐步学习如何把边缘、深度、分割图、人体姿态等条件稳定注入到生成过程中。其优点在于既保留了大模型原有的生成能力，又能以较小成本增加外部结构约束，从而显著提升可控生成能力。

对于视频风格生成任务而言，ControlNet的意义主要体现在两个方面。第一，它能够在角色姿态、场景布局和镜头结构等方面提供更明确的条件约束，从而降低生成过程中的构图漂移；第二，它与LoRA之间具有天然的互补关系，LoRA更偏向纹理风格与视觉属性的学习，而ControlNet更偏向几何结构和空间布局的保持\cite{ref14}。因此，在既要保证画风接近目标风格、又要尽量维持主体结构稳定的场景中，二者联合使用往往比单独依赖其中任意一种方法更有效。

\subsection{风格保持与主体一致性问题}

在特定风格视频生成任务中，风格保持与主体一致性是两个相互关联、又彼此制约的目标。若过度强调风格迁移，模型可能会牺牲主体结构、动作合理性和跨帧稳定性；若过度强调结构一致性，则可能导致风格表达不足，画面趋于保守。AnimateDiff\cite{ref12}和ControlNet\cite{ref14}等工作分别从时间运动建模与空间结构控制两个角度对这一问题提供了解法。结合LoRA的参数高效微调能力，可以形成一种更符合工程实践的组合式思路：以LoRA学习外观风格，以运动模块保证时序连贯，以结构控制模块维持主体姿态和场景布局。

本研究的技术路线正是在这一理论认识基础上逐步展开的。前期先使用Stable Diffusion底模与LoRA进行图像风格学习，验证特定风格是否能够稳定注入；中期再通过批量推理、图像对比和可视化分析，评估LoRA对色彩、构图和风格强度的影响；后期则进一步扩展到视频阶段，考察在加入时序建模与结构控制后，能否在保持画风相似性的同时获得更平滑的动作表现和更稳定的主体特征。由此可以看出，风格保持与一致性控制并不是彼此孤立的模块，而是构成特定风格视频生成系统的两条主线。

\section{本章小结}

本章围绕本研究涉及的核心理论与关键技术进行了系统梳理。首先介绍了扩散模型的基本原理以及条件引导机制，明确了当前文生图任务的概率建模基础；其次分析了潜空间扩散模型和Stable Diffusion的网络结构，说明了其在高分辨率生成与可控条件注入方面的优势；随后重点阐述了LoRA参数高效微调技术及其在扩散模型中的应用方式，解释了本研究为何能够在有限算力条件下完成风格学习实验；最后，从CLIP/BLIP数据标注、视频扩散、AnimateDiff以及ControlNet等角度，说明了本研究在数据构建、视频扩展与结构控制方面所依赖的关键技术。
